% Shared preamble for the Integration training sheet.
\usepackage[T1]{fontenc}
\usepackage[utf8]{inputenc}
\usepackage{lmodern}
\ifdefined\KindleEdition
  % 6-inch e-ink page (Kindle Paperwhite): the page *is* the screen, so the
  % margins are only what the eye needs, not what a printer needs.
  \usepackage[paperwidth=90mm,paperheight=122mm,
              top=4mm,bottom=6mm,left=4.5mm,right=4.5mm]{geometry}
\else
  \usepackage[a4paper,margin=2.4cm]{geometry}
\fi
\usepackage{amsmath,amssymb,amsthm}
\usepackage{microtype}
\usepackage{enumitem}
\usepackage{xcolor}
\usepackage[most]{tcolorbox}
\usepackage{fancyhdr}
\usepackage[hidelinks]{hyperref}

\definecolor{accent}{HTML}{1F4E79}
\definecolor{soft}{HTML}{F4F7FA}

% --- function graphs ------------------------------------------------------
\usepackage{pgfplots}
\pgfplotsset{compat=1.18}
\usetikzlibrary{arrows.meta}

% One consistent look for every function graph in these books.
\pgfplotsset{
  fnaxis/.style={
    width=5.9cm, height=4.6cm,
    axis lines=middle,
    axis line style={-{Stealth[length=4pt]}, gray!65!black, line width=0.3pt},
    tick style={gray!65!black},
    tick label style={font=\tiny, gray!50!black},
    label style={font=\scriptsize},
    xlabel={$x$}, ylabel={$y$},
    every axis x label/.style={at={(ticklabel* cs:1.0)}, anchor=west},
    every axis y label/.style={at={(ticklabel* cs:1.0)}, anchor=south east,
      xshift=-1pt},
    every axis plot/.append style={accent, very thick, smooth, samples=201},
    clip=true, enlargelimits=false,
    title style={font=\scriptsize, yshift=4pt},
  },
  fnaxis/.append style={/pgfplots/legend style={font=\tiny, draw=none,
    fill=none, at={(0.98,0.03)}, anchor=south east}},
}

% A row of small graphs, evenly spaced with \hfill between them, plus an
% optional caption line: \begin{gallery} <pic>\hfill<pic> \gcap{...} \end{gallery}
\newenvironment{gallery}
  {\par\medskip\noindent\ignorespaces}
  {\par\medskip\noindent}
\newcommand{\gcap}[1]{\par\smallskip{\centering\footnotesize #1\par}}

% --- plane regions, implicit curves, surfaces -----------------------------
\usepgfplotslibrary{fillbetween}
\definecolor{regionfill}{HTML}{9DB4D8}
\pgfplotsset{colormap={house}{color=(soft) color=(accent!45) color=(accent)}}

\pgfplotsset{
  % Squared, gridded axis for plane regions and implicitly given curves.
  geoaxis/.style={
    fnaxis,
    width=5.9cm, height=5.9cm,
    grid=both,
    grid style={gray!22, line width=0.2pt},
    minor tick num=1,
  },
  % Level-curve maps: many thin curves, no single curve emphasised.
  contouraxis/.style={
    geoaxis,
    every axis plot/.append style={accent!75, line width=0.5pt, smooth,
      samples=201},
  },
  % Surfaces: \addplot3[surf, domain=..., y domain=...] {expr};
  surfaxis/.style={
    width=6.0cm, height=5.0cm,
    view={130}{22},
    colormap name=house,
    ticks=none,
    xlabel={$x$}, ylabel={$y$}, zlabel={$f$},
    label style={font=\tiny, gray!50!black},
    title style={font=\scriptsize},
    every axis plot/.append style={samples=25, shader=faceted interp,
      draw=accent!35, line width=0.1pt},
  },
}

% Shorthand for a shaded region between two named paths.
% \region{<path A>}{<path B>}{<soft clip options>}
\newcommand{\region}[3]{%
  \addplot[regionfill, opacity=0.55, forget plot, sharp plot]
    fill between[of=#1 and #2, #3];}

% --- rule boxes -----------------------------------------------------------
\newtcolorbox{rulebox}[1]{%
  breakable, enhanced, colback=soft, colframe=accent,
  boxrule=0.4pt, left=7pt, right=7pt, top=6pt, bottom=6pt,
  fonttitle=\bfseries, coltitle=white, colbacktitle=accent,
  title={#1}}

% --- worked examples ------------------------------------------------------
% One fully narrated solution per chapter, carrying the chapter's central
% technique and the mistakes that technique invites:
%   \begin{worked}{Task} ... \wstep{Step 1.} ... \wtrap{...} ... \end{worked}
\newtcolorbox{worked}[1]{%
  breakable, enhanced, colback=white, colframe=accent!65,
  boxrule=0.4pt, left=7pt, right=7pt, top=6pt, bottom=6pt,
  fonttitle=\bfseries, coltitle=accent, colbacktitle=soft,
  title={#1}}

% Run-in step heading inside a worked example.
\newcommand{\wstep}[1]{\par\smallskip\noindent\textbf{#1}\enspace\ignorespaces}

% A pitfall, set off from the running solution so it can be found again.
\newcommand{\wtrap}[1]{%
  \par\smallskip
  {\leftskip=1.3em\noindent
   \llap{\textcolor{accent}{\rule[0.22ex]{0.9em}{0.6pt}}\hspace{0.4em}}%
   \textbf{Trap.}\enspace\itshape #1\par}%
  \smallskip}

% --- problem / solution lists --------------------------------------------
\newlist{problems}{enumerate}{1}
\setlist[problems]{label=\textbf{\arabic*.}, leftmargin=2.1em,
  itemsep=0.55ex, parsep=0pt, topsep=0.8ex}
\newlist{solutions}{enumerate}{1}
\setlist[solutions]{label=\textbf{\arabic*.}, leftmargin=2.1em,
  itemsep=1.0ex, parsep=0pt, topsep=0.8ex}

% --- running heads --------------------------------------------------------
\pagestyle{fancy}
\fancyhf{}
\fancyhead[LE,RO]{\thepage}
\fancyhead[RE]{\small\itshape\nouppercase{\leftmark}}
\fancyhead[LO]{\small\itshape\nouppercase{\rightmark}}
\setlength{\headheight}{14pt}
\renewcommand{\headrulewidth}{0.3pt}
\renewcommand{\chaptermark}[1]{\markboth{\chaptername\ \thechapter.\ #1}{}}
\renewcommand{\sectionmark}[1]{\markright{\thesection\ #1}}

% --- shorthands used throughout -------------------------------------------
\newcommand{\dd}[2]{\frac{\mathrm{d}#1}{\mathrm{d}#2}}
\newcommand{\pd}[2]{\frac{\partial #1}{\partial #2}}
\newcommand{\pdd}[3]{\frac{\partial^{2} #1}{\partial #2\,\partial #3}}
\newcommand{\R}{\mathbb{R}}
\newcommand{\grad}{\nabla}
\DeclareMathOperator{\sgn}{sgn}
\DeclareMathOperator{\Hess}{H}

% German gloss, used on first mention of a technical term
\newcommand{\de}[1]{\emph{#1}}

% --- cross-references to chapters ------------------------------------------
% Chapter files are shared by the standalone volumes and the compendium, so a
% hard-coded "Chapter 6" is wrong in one of them. Write \chref{derivatives}{6}
% instead: it prints "Chapter 6" inside the same volume, "Volume 5, Chapter 6"
% from another volume, and the running number ("Chapter 34") in the
% compendium. \chrefs{vol}{n} is the abbreviated form (Ch. / Vol., Ch.),
% \chnum{vol}{n} the bare number in the current build, and \volref{vol} names
% the volume ("Volume 5", or "Part V" in the compendium). Each volume's
% main.tex sets \thisvolume; the compendium's sets \incompendiumtrue.
\newif\ifincompendium
\newcommand*{\thisvolume}{}
\makeatletter
\@namedef{chref@vol@algebra}{1}\@namedef{chref@off@algebra}{0}
\@namedef{chref@vol@functions}{2}\@namedef{chref@off@functions}{9}
\@namedef{chref@vol@matrices}{3}\@namedef{chref@off@matrices}{14}
\@namedef{chref@vol@foundations}{4}\@namedef{chref@off@foundations}{24}
\@namedef{chref@vol@derivatives}{5}\@namedef{chref@off@derivatives}{28}
\@namedef{chref@vol@integrals}{6}\@namedef{chref@off@integrals}{42}
\@namedef{chref@vol@optimization}{7}\@namedef{chref@off@optimization}{50}
\@namedef{chref@vol@linprog}{8}\@namedef{chref@off@linprog}{60}
\@namedef{chref@vol@mixed}{9}\@namedef{chref@off@mixed}{68}
\newcommand*{\chref@check}[1]{%
  \@ifundefined{chref@vol@#1}{\PackageError{chref}{Unknown volume `#1'}{}}{}}
\newcommand*{\chnum}[2]{\chref@check{#1}%
  \ifincompendium\the\numexpr\csname chref@off@#1\endcsname+#2\relax
  \else#2\fi}
% \chref@same{vol}{same volume}{other volume}
\newcommand*{\chref@same}[3]{\def\chref@tmp{#1}%
  \ifx\chref@tmp\thisvolume\expandafter\@firstoftwo
  \else\expandafter\@secondoftwo\fi{#2}{#3}}
\newcommand*{\chref}[2]{\chref@check{#1}%
  \ifincompendium Chapter~\chnum{#1}{#2}%
  \else\chref@same{#1}{Chapter~#2}%
    {Volume~\csname chref@vol@#1\endcsname, Chapter~#2}\fi}
\newcommand*{\chrefs}[2]{\chref@check{#1}%
  \ifincompendium Ch.~\chnum{#1}{#2}%
  \else\chref@same{#1}{Ch.~#2}%
    {Vol.~\csname chref@vol@#1\endcsname, Ch.~#2}\fi}
\newcommand*{\volref}[1]{\chref@check{#1}%
  \ifincompendium Part~\@Roman{\csname chref@vol@#1\endcsname}%
  \else Volume~\csname chref@vol@#1\endcsname\fi}
\makeatother

% --- Kindle edition: compact layout for a 6-inch screen -------------------
% Everything below applies only when \KindleEdition is defined on the
% command line; the A4 edition is untouched.
\ifdefined\KindleEdition
  \usepackage{etoolbox}
  \usepackage{adjustbox}
  \usepackage{titlesec}

  % The 8cm measure leaves few break points; let TeX loosen a tight line
  % rather than push it into the margin.
  \setlength{\emergencystretch}{1.5em}

  % No running heads: on a 90mm page the header costs more than it gives.
  % A centred page number in the foot is kept so the contents list means
  % something.
  \pagestyle{plain}
  \fancypagestyle{plain}{\fancyhf{}\fancyfoot[C]{\small\thepage}%
    \renewcommand{\headrulewidth}{0pt}}

  % Compact headings: the book class leaves a third of a page above a
  % chapter title, which is unaffordable here.
  \titleformat{\chapter}[display]
    {\normalfont\Large\bfseries\color{accent}}{\normalsize\chaptertitlename\ \thechapter}{0.4em}{\Large}
  \titlespacing*{\chapter}{0pt}{0pt}{1.2em}
  \titlespacing*{\section}{0pt}{1.4ex plus .3ex}{0.7ex}
  \titlespacing*{\subsection}{0pt}{1.2ex plus .2ex}{0.5ex}

  % Boxes: trim the padding, which was set for a 16cm measure.
  \tcbset{left=3.5pt,right=3.5pt,top=3pt,bottom=3pt,
          fonttitle=\bfseries\small,
          halign title=flush left}

  % Problem and solution lists: reclaim the wide left margin.
  \setlist[problems]{label=\textbf{\arabic*.}, leftmargin=1.5em,
    itemsep=0.45ex, parsep=0pt, topsep=0.6ex}
  \setlist[solutions]{label=\textbf{\arabic*.}, leftmargin=1.5em,
    itemsep=0.8ex, parsep=0pt, topsep=0.6ex}

  % Figures and wide tables shrink to the measure instead of running off it.
  % adjustbox only scales down: anything already narrow is left alone.
  \BeforeBeginEnvironment{tikzpicture}{\begin{adjustbox}{max width=\linewidth}}
  \AfterEndEnvironment{tikzpicture}{\end{adjustbox}}
  \BeforeBeginEnvironment{tabular}{\begin{adjustbox}{max width=\linewidth}}
  \AfterEndEnvironment{tabular}{\end{adjustbox}}

  % Display maths was written for a 16cm measure. On 8cm the long chained
  % equations would run off the screen, so any display wider than the text
  % block is scaled down to fit; narrower ones are left at full size. The
  % whole corpus uses \[ ... \] and no other display environment, so this
  % one redefinition covers every case.
  \newsavebox{\kndbox}
  \renewcommand{\[}{\par\addvspace{0.7ex}%
    \begin{lrbox}{\kndbox}$\displaystyle}
  \renewcommand{\]}{$\end{lrbox}%
    {\centering
      \ifdim\wd\kndbox>\linewidth
        \resizebox{\linewidth}{!}{\usebox{\kndbox}}%
      \else
        \usebox{\kndbox}%
      \fi
     \par}%
    \addvspace{0.7ex}}

  % A narrow measure makes perfect justification impossible; prefer a loose
  % line to a line that runs into the margin.
  \sloppy
  \raggedbottom
\fi
